\subsection{Casos de Uso de la Defensoría}
\label{subsec:casos_uso_defensoria}

La figura \ref{fig:casos_uso_defensoria_titular} ilustra las funciones del Titular de la Defensoría, quien realiza la supervisión final de los expedientes y gestiona el cierre administrativo de los casos resueltos.

\begin{figure}[H]
    \centering
    \includegraphics[width=0.7\textwidth]{imagenes/defensoria/Diagrama_CU/CasosUso_Defensoria.png}
    \caption{Diagrama de Casos de Uso del Microservicio Defensoría}
    \label{fig:casos_uso_defensoria_titular}
\end{figure}

\vspace{0.5cm}

% ================= CU65 =================
\begingroup
\setlength{\parskip}{0pt}
\begin{longtable}{|p{3.5cm}|p{2cm}|p{2.5cm}|p{7cm}|}
\caption{Especificación de Caso de Uso CU65: Consultar panel de defensoría}
\label{tab:cu65} \\
\hline
\textbf{Identificador} & CU65 & \textbf{Nombre} & Consultar panel de defensoría \\ \hline
\textbf{Actor Principal} & \multicolumn{3}{p{11.5cm}|}{Defensoría} \\ \hline
\textbf{Descripción} & \multicolumn{3}{p{11.5cm}|}{Acceso a la bandeja global de quejas que han finalizado su proceso jurídico y requieren cierre.} \\ \hline
\textbf{Disparador} & \multicolumn{3}{p{11.5cm}|}{El usuario accede al módulo de alta dirección.} \\ \hline
\textbf{Precondiciones} & \multicolumn{3}{p{11.5cm}|}{Sesión activa con rol de Defensor Titular.} \\ \hline
\textbf{Flujo Principal} & \multicolumn{3}{p{11.5cm}|}{1. El Sistema consulta expedientes en estatus 'FIRMADO' o 'POR CERRAR'.\newline 2. El Sistema muestra una tabla resumen con folios, fechas y resoluciones.} \\ \hline
\textbf{Flujos Alternativos} & \multicolumn{3}{p{11.5cm}|}{Búsqueda por filtros avanzados de periodos anuales.} \\ \hline
\textbf{Situaciones de error} & \multicolumn{3}{p{11.5cm}|}{Error de carga por saturación de base de datos.} \\ \hline
\textbf{Estado en error} & \multicolumn{3}{p{11.5cm}|}{El sistema muestra mensaje de espera.} \\ \hline
\textbf{Postcondiciones} & \multicolumn{3}{p{11.5cm}|}{Lista de expedientes lista para supervisión.} \\ \hline
\end{longtable}
\endgroup

\vspace{0.8cm}

% ================= CU66 =================
\begingroup
\setlength{\parskip}{0pt}
\begin{longtable}{|p{3.5cm}|p{2cm}|p{2.5cm}|p{7cm}|}
\caption{Especificación de Caso de Uso CU66: Abrir queja para cierre}
\label{tab:cu66} \\
\hline
\textbf{Identificador} & CU66 & \textbf{Nombre} & Abrir queja para cierre \\ \hline
\textbf{Actor Principal} & \multicolumn{3}{p{11.5cm}|}{Defensoría} \\ \hline
\textbf{Descripción} & \multicolumn{3}{p{11.5cm}|}{Visualización del expediente completo antes de proceder al archivo definitivo.} \\ \hline
\textbf{Disparador} & \multicolumn{3}{p{11.5cm}|}{Selección de un folio en el panel (CU65).} \\ \hline
\textbf{Precondiciones} & \multicolumn{3}{p{11.5cm}|}{El expediente debe estar debidamente firmado por la Subdefensoría.} \\ \hline
\textbf{Flujo Principal} & \multicolumn{3}{p{11.5cm}|}{1. El Sistema despliega la información del quejoso, hechos y la resolución final.\newline 2. Permite la descarga del documento firmado para revisión visual.} \\ \hline
\textbf{Flujos Alternativos} & \multicolumn{3}{p{11.5cm}|}{N/A.} \\ \hline
\textbf{Situaciones de error} & \multicolumn{3}{p{11.5cm}|}{Documento final no encontrado en el servidor de archivos.} \\ \hline
\textbf{Estado en error} & \multicolumn{3}{p{11.5cm}|}{Aviso de error técnico; se impide el cierre hasta que el archivo sea restaurado.} \\ \hline
\textbf{Postcondiciones} & \multicolumn{3}{p{11.5cm}|}{Expediente listo para confirmación de cierre.} \\ \hline
\end{longtable}
\endgroup

\vspace{0.8cm}

% ================= CU67 =================
\begingroup
\setlength{\parskip}{0pt}
\begin{longtable}{|p{3.5cm}|p{2cm}|p{2.5cm}|p{7cm}|}
\caption{Especificación de Caso de Uso CU67: Cerrar queja definitivamente}
\label{tab:cu67} \\
\hline
\textbf{Identificador} & CU67 & \textbf{Nombre} & Cerrar queja definitivamente \\ \hline
\textbf{Actor Principal} & \multicolumn{3}{p{11.5cm}|}{Defensoría} \\ \hline
\textbf{Descripción} & \multicolumn{3}{p{11.5cm}|}{Acción formal que concluye el trámite administrativo de la queja.} \\ \hline
\textbf{Disparador} & \multicolumn{3}{p{11.5cm}|}{El Defensor selecciona la opción "Cerrar Expediente".} \\ \hline
\textbf{Precondiciones} & \multicolumn{3}{p{11.5cm}|}{Resolución enviada y notificada previamente.} \\ \hline
\textbf{Flujo Principal} & \multicolumn{3}{p{11.5cm}|}{1. El usuario confirma el cierre.\newline 2. El Sistema incluye CU68 (Cambiar estatus a CONCLUIDA).\newline 3. El Sistema incluye CU69 (Notificar cierre).} \\ \hline
\textbf{Flujos Alternativos} & \multicolumn{3}{p{11.5cm}|}{N/A.} \\ \hline
\textbf{Situaciones de error} & \multicolumn{3}{p{11.5cm}|}{Fallo en la actualización del registro en el microservicio de expedientes.} \\ \hline
\textbf{Estado en error} & \multicolumn{3}{p{11.5cm}|}{El sistema revierte el cambio de estado (Rollback).} \\ \hline
\textbf{Postcondiciones} & \multicolumn{3}{p{11.5cm}|}{La queja pasa a estatus de 'Concluida'.} \\ \hline
\end{longtable}
\endgroup

\vspace{0.8cm}

% ================= CU68 =================
\begingroup
\setlength{\parskip}{0pt}
\begin{longtable}{|p{3.5cm}|p{2cm}|p{2.5cm}|p{7cm}|}
\caption{Especificación de Caso de Uso CU68: Cambiar estatus a concluida}
\label{tab:cu68} \\
\hline
\textbf{Identificador} & CU68 & \textbf{Nombre} & Cambiar estatus a concluida \\ \hline
\textbf{Actor Principal} & \multicolumn{3}{p{11.5cm}|}{Sistema} \\ \hline
\textbf{Descripción} & \multicolumn{3}{p{11.5cm}|}{Proceso automatizado para marcar el folio como terminado en la base de datos.} \\ \hline
\textbf{Disparador} & \multicolumn{3}{p{11.5cm}|}{Confirmación de cierre en CU67.} \\ \hline
\textbf{Precondiciones} & \multicolumn{3}{p{11.5cm}|}{Folio activo en fase final.} \\ \hline
\textbf{Flujo Principal} & \multicolumn{3}{p{11.5cm}|}{1. El Sistema cambia el atributo estatus a 'CONCLUIDA'.\newline 2. El Sistema bloquea cualquier modificación posterior al expediente.} \\ \hline
\textbf{Flujos Alternativos} & \multicolumn{3}{p{11.5cm}|}{N/A.} \\ \hline
\textbf{Situaciones de error} & \multicolumn{3}{p{11.5cm}|}{Pérdida de conexión con el motor de base de datos.} \\ \hline
\textbf{Estado en error} & \multicolumn{3}{p{11.5cm}|}{Registro de error en logs del sistema.} \\ \hline
\textbf{Postcondiciones} & \multicolumn{3}{p{11.5cm}|}{Expediente archivado electrónicamente.} \\ \hline
\end{longtable}
\endgroup

\vspace{0.8cm}

% ================= CU69 =================
\begingroup
\setlength{\parskip}{0pt}
\begin{longtable}{|p{3.5cm}|p{2cm}|p{2.5cm}|p{7cm}|}
\caption{Especificación de Caso de Uso CU69: Notificar cierre de queja}
\label{tab:cu69} \\
\hline
\textbf{Identificador} & CU69 & \textbf{Nombre} & Notificar cierre de queja \\ \hline
\textbf{Actor Principal} & \multicolumn{3}{p{11.5cm}|}{Sistema} \\ \hline
\textbf{Descripción} & \multicolumn{3}{p{11.5cm}|}{Envío de un correo final de cortesía informando que el expediente ha sido archivado.} \\ \hline
\textbf{Disparador} & \multicolumn{3}{p{11.5cm}|}{Éxito en el cambio de estatus (CU68).} \\ \hline
\textbf{Precondiciones} & \multicolumn{3}{p{11.5cm}|}{Correo electrónico del quejoso disponible.} \\ \hline
\textbf{Flujo Principal} & \multicolumn{3}{p{11.5cm}|}{1. El Sistema redacta el mensaje de conclusión de trámite.\newline 2. Se envía el correo al usuario.\newline 3. Se registra el envío en el historial del folio.} \\ \hline
\textbf{Flujos Alternativos} & \multicolumn{3}{p{11.5cm}|}{N/A.} \\ \hline
\textbf{Situaciones de error} & \multicolumn{3}{p{11.5cm}|}{Fallo del servidor SMTP.} \\ \hline
\textbf{Estado en error} & \multicolumn{3}{p{11.5cm}|}{El sistema marca la notificación como fallida para revisión manual.} \\ \hline
\textbf{Postcondiciones} & \multicolumn{3}{p{11.5cm}|}{El quejoso es informado de la conclusión administrativa.} \\ \hline
\end{longtable}
\endgroup